\documentclass[11pt]{article}
\usepackage[a4paper,margin=25mm]{geometry}
\usepackage[T1]{fontenc}
\usepackage{lmodern}
\usepackage{booktabs}
\usepackage{amsmath}
\usepackage{graphicx}
\usepackage{hyperref}
\usepackage{xcolor}
\usepackage{listings}
\lstset{basicstyle=\ttfamily\small, columns=fullflexible, keepspaces=true,
  breaklines=true, frame=single, framerule=0.3pt, rulecolor=\color{gray}}
\newcommand{\code}[1]{\texttt{#1}}
\newcommand{\us}{\textmu s}

\title{\code{template.py}: the shape of the single-photoelectron pulse}
\author{BGOTSCAN}
\date{2026-10-14}

\begin{document}
\maketitle

\section{Purpose}

The single-photoelectron pulses of the SiPMs are 15 counts high and two or
three samples wide (about 20~ns) on 4 counts of white noise per sample. A
threshold on one sample cannot separate them from noise
(\code{spe\_doc.tex}): a matched filter, which uses the shape of the pulse,
is needed, and the shape must be measured first. This script does only
that: it finds the pulse shape, the \emph{template}, from the pulses of
the records, checks it channel by channel, and writes it for
\code{spe.py} to use. It processes one event at a time in six steps, with
the text and one figure per step, as \code{selection.py} and
\code{spe.py}; the shape needs thousands of pulses, so its normal use is
the \code{--summary} mode over a few hundred events. It is a Python script
in \code{src/py/}, reading the raw HDF5 files with h5py.

\section{Usage}

\begin{lstlisting}
python3 src/py/template.py RUN [options]
\end{lstlisting}

Interactive: after each step the script waits on the terminal; Enter goes
to the next step, \code{n} abandons the event and starts the next one,
\code{q} quits.

\begin{center}
\begin{tabular}{l l p{74mm}}
\toprule
option & default & meaning \\
\midrule
\code{--data DIR} & the data directory & as in \code{wfview.py} \\
\code{--events SPEC} & \code{0} & events to process: \code{a-b} or \code{a,b,c}, 0-based \\
\code{--channel C} & 0 & the channel shown in steps 1 and 2 \\
\code{--region R} & \code{tail} & search region: \code{tail}, the last \code{--tail} \us{} of the record, or \code{all} \\
\code{--tail T} & 10 & length of the tail region in \us \\
\code{--threshold A} & 12 & candidate threshold on the raw waveform, in ADC counts (3 sigma: high on purpose) \\
\code{--half H} & 8 & half length of the template in samples: $2H + 1$ samples, 136~ns for $H = 8$ \\
\code{--from-step N} & 1 & start the display at this step \\
\code{--save NAME} & off & no prompts; figures saved as \code{src/py/plots/NAME\_evE\_sN.png} \\
\code{--summary} & off & steps 1 and 2 silent per event; steps 3 to 6 once after the last event (figures saved as \code{NAME\_sN.png} with \code{--save}) \\
\code{--write FILE} & off & write the refined template of step~5 to \code{FILE} as CSV \\
\bottomrule
\end{tabular}
\end{center}

\section{Notation}

$x_i$ is a channel waveform minus its baseline, the median over the
record; $\Delta t = 8$~ns; $R$ is the search region (\code{spe\_doc.tex},
step~1). No common-mode subtraction. A template is a vector $T_j$, $j =
-H \dots H$, normalised to $T_0 = 1$ at the peak.

\section{The steps}

\subsection{Step 1: candidates on the channel}

Candidates are the local maxima of the raw $x$ inside $R$ with $x >
A_{\rm thr}$ (\code{--threshold}, 12 counts), taken in decreasing order
and at least $2H$ samples apart. The threshold is high so that most
candidates are real pulses; the price, pulses lost on the low side, does
not matter here, since only the shape is wanted. Text: number, times and
amplitudes. Figure: $x$ over $R$ with the threshold and the candidates
marked.

\subsection{Step 2: aligned pulses of the channel}

For each candidate at sample $i_A$ (its maximum), the raw samples $x_{i_A
- H} \dots x_{i_A + H}$ minus a local baseline, the median of the window
samples more than two samples before the peak, which removes the slow
wander of the channel at the time of the pulse without being pulled by
the tail of the pulse itself. Text: the number of
pulses. Figure: the pulses
overlaid against the sample offset $j$, and their mean.

\subsection{Step 3: all channels, accumulated}

Steps 1 and 2 on the 64 channels, added to the pulses accumulated over
the events processed so far. Text: pulses per channel as an $8\times8$
table and the total. Figure: a two-dimensional histogram of amplitude
versus offset $j$ over all accumulated pulses, with their mean drawn on
top: the shape with its spread.

\subsection{Step 4: first template}

$T^{(1)}_j$ is the mean of the accumulated pulses at offset $j$, divided
by its value at $j = 0$. Text: the $2H + 1$ values, the full width at
half maximum in samples (linear interpolation between samples), the
undershoot ($\min_{j > 0} T_j$), the sum of squares $\sum_j T_j^2$, and
the standard error of the mean at the peak. Figure: $T^{(1)}$ with the
standard error of each sample.

The sum of squares is what the matched filter gains. For a pulse of
shape $T$ on white noise $\sigma$, the least-squares amplitude at sample
$i$ is
\[
  y_i = \frac{\sum_j T_j\, x_{i+j}}{\sum_j T_j^2},
  \qquad \sigma_y = \frac{\sigma}{\sqrt{\sum_j T_j^2}},
\]
so the signal to noise improves by $\sqrt{\sum_j T_j^2}$ over the single
sample.

Bias: the candidates were selected on their maximum sample, so $x_0$ is
on average high by the upward noise that helped it pass the threshold,
about one count on 15, and the neighbours are not. $T^{(1)}$ is therefore
slightly too peaked. The next step removes most of it.

\subsection{Step 5: refined template}

Every accumulated channel waveform is filtered with $T^{(1)}$ (the $y_i$
above); candidates are found on $y$ with a threshold of $A_{\rm thr}$
expressed in units of $\sigma_y$, measured as the rms of $y$ over $R$
after clipping at 3 sigma; the raw samples are aligned on the maximum of
$y$ and averaged: $T^{(2)}$. Since the alignment is now on a filtered
quantity, the selection does not favour the central raw sample. Text:
the values of $T^{(2)}$, the change $T^{(2)}_j - T^{(1)}_j$ per sample,
both sums of squares. Figure: both templates overlaid. $T^{(2)}$ is the
template to use.

\subsection{Step 6: per-channel check}

For each channel, the mean of its own aligned pulses (from step~5),
normalised to its peak, against $T^{(2)}$: the residual $r_c$ is the rms
over the $2H + 1$ samples of the difference. Text: the median of $r_c$,
and the channels with $r_c$ above three times the median. Figure: the
$8\times8$ map of $r_c$. If channels stand out, the pulse shape depends
on the channel and the pass needs one template per channel; otherwise
one template serves all.

\section{Output}

With \code{--write FILE}, $T^{(2)}$ is written as a CSV with two columns,
offset in samples and value, and a header giving the run, the number of
pulses, the threshold and $\sum_j T_j^2$. The natural place is
\code{src/spe\_template\_RUN.csv}, tracked; whether one template serves all
runs is decided by comparing them.

\section{Examples}

\begin{lstlisting}
# one event, channel 0, every step (the template from a handful of pulses)
python3 src/py/template.py 197K_39V_th20

# the template from 300 events, figures saved, template written
python3 src/py/template.py 197K_39V_th20 --events 0-299 --summary --save tpl197 \
        --write src/spe_template_197K_39V_th20.csv
\end{lstlisting}

\end{document}
